%!TEX TS-program = pdflatex
\documentclass[aspectratio=169,10pt]{beamer}
\usepackage[T2A]{fontenc}
\usepackage[utf8]{inputenc}
\usepackage[english,russian]{babel}
\usepackage{cmap}
\usepackage{amsmath,amssymb,mathtools}
\usepackage{booktabs,array}
\usepackage{etoolbox}
\usepackage{PTSans}
\usepackage{graphicx}
\graphicspath{{figures/}}
\renewcommand{\familydefault}{\sfdefault}
\renewcommand{\rmdefault}{\sfdefault}
\newbool{russian}\booltrue{russian}
\usepackage{HSE-theme/beamerthemeHSE}
\hypersetup{bookmarks=false}
\setbeamertemplate{blocks}[rounded][shadow=false]
\setbeamerfont{frametitle}{size=\Large}
\setbeamerfont{framesubtitle}{size=\small}
\setbeamerfont{block title}{size=\normalsize}
\setbeamerfont{block body}{size=\small}
\newcommand{\R}{\mathbb{R}}
\newcommand{\norm}[1]{\left\lVert #1\right\rVert}
\newcommand{\grad}{\nabla}
\newcommand{\jac}{J}
\newcommand{\pic}[2][.83]{\includegraphics[width=#1\linewidth]{#2}}
\title[Методы оптимизации в ML]{Методы оптимизации\\в машинном обучении}
\subtitle{Лекция 7. Матричное и автоматическое дифференцирование}
\author[Н. Лукьяненко]{Никита Лукьяненко}
\institute{Совместная программа СМОЛ ГУ и ФКН НИУ ВШЭ}
\date{2026/27 учебный год}
\begin{document}

% 01
\frame[plain]{\titlepage}

% 02
\begin{frame}{Градиентный спуск: откуда берётся градиент?}
\small
За шесть лекций мы научились исследовать метод
\[
\boxed{x_{k+1}=x_k-\alpha_k\grad f(x_k).}
\]
Мы обсуждали кривизну, число обусловленности и скорость сходимости.

\begin{alertblock}{Скрытая вычислительная задача}
Каждая итерация требует вычислить $\grad f(x_k)$. Что делать, если $f$ включает тысячи параметров и промежуточных операций?
\end{alertblock}
\end{frame}

% 03
\begin{frame}{Функцию просто вычислить, но сложно дифференцировать}
\small
Для $f(x)=x^2+\sin x$ производная очевидна. Но даже композиция
\[
 f(x)=\bigl(\sin(x^2+e^x)+x^3\bigr)^2
\]
выглядит намного менее дружелюбно после полного раскрытия производной.

\centering
\pic[.84]{long_composition.pdf}

\small
\vspace{.2em}
\textbf{Идея:} дифференцировать элементарные операции по отдельности.
\end{frame}

% 04
\begin{frame}{Автоматическое дифференцирование: точное определение}
\footnotesize
\begin{block}{Определение}
\textbf{Автоматическое дифференцирование} (АД) — это вычисление производных функции, заданной последовательностью элементарных операций, посредством систематического применения \textbf{правила цепочки} к каждой операции.
\end{block}
\vspace{.25em}
\begin{itemize}
\item \textbf{Символьный подход:} строим формулу производной как математическое выражение.
\item \textbf{Численный подход:} приближаем производную конечными разностями с шагом $h$.
\item \textbf{Автоматический подход:} на конкретных входных данных вычисляем значения операций и их производные по точным аналитическим правилам.
\end{itemize}
\begin{alertblock}{Важно}
АД не использует конечные разности и не требует раскрывать полную символическую формулу. Точность всё равно ограничена машинной арифметикой.
\end{alertblock}
\end{frame}
% 05
\begin{frame}{Карта сегодняшней лекции}
\small
\[
\boxed{\text{Дифференциал}}
\Longrightarrow
\boxed{\text{Правило цепочки}}
\]
\[
\boxed{\text{Вычислительный граф}}
\Longrightarrow
\boxed{\text{Автоматический градиент}}.
\]

В конце лекции самостоятельно пройдём по графу функции
\[
\boxed{f(x,y)=(xy+\sin x)^2}
\]
и получим производные в двух разных режимах.
\end{frame}

% 06
\begin{frame}{Возвращаемся к линейному приближению}
\small
Из лекции №2 для дифференцируемой $f:\R^d\to\R$:
\[
f(x+h)=f(x)+\grad f(x)^\top h+o(\norm{h}).
\]
\centering
\pic[.73]{differential_tangent.pdf}
\end{frame}

% 07
\begin{frame}{Дифференциал — главная линейная часть изменения}
\small
Приращение $dx$ может быть малым вектором произвольного направления. Определяем дифференциал как линейную по $dx$ часть приращения:
\[
\boxed{df=\grad f(x)^\top dx.}
\]
Важно различать:
\[
 f(x+dx)-f(x)=df+o(\norm{dx}).
\]
\begin{alertblock}{Не путать}
$df$ не обязано совпадать с настоящим конечным приращением функции. Совпадает лишь его линейная часть.
\end{alertblock}
\end{frame}

% 08
\begin{frame}{Пример: как прочитать градиент из дифференциала}
\small
Пусть $f(x,y)=x^2+3xy$. Тогда
\[
 df=(2x+3y)\,dx+3x\,dy.
\]
Соберём коэффициенты в столбец:
\[
 df=
\underbrace{\begin{pmatrix}2x+3y&3x\end{pmatrix}}_{\grad f(x,y)^\top}
\begin{pmatrix}dx\\dy\end{pmatrix}.
\]
\begin{block}{Рабочий приём}
Привести $df$ к форме $g^\top dx$, после чего прочитать $\grad f=g$.
\end{block}
\end{frame}

% 09
\begin{frame}{А если результат — целый вектор?}
\small
Пусть $g:\R^d\to\R^m$, тогда дифференциал тоже является вектором:
\[
 dg=J_g(x)\,dx,
\]
где $J_g(x)$ — матрица линейного приближения.

\[
\underbrace{dg}_{m\times1}
=\underbrace{J_g(x)}_{m\times d}
\underbrace{dx}_{d\times1}.
\]
Определение матрицы Якоби подробно появится чуть позже.
\end{frame}

% 10
\begin{frame}{Проверка размерностей — лучший способ найти ошибку}
\centering
\pic[.82]{matrix_dimensions.pdf}

\small
Для любого дифференцируемого отображения $g:\R^d\to\R^m$ размерность якобиана определяется автоматически: $m\times d$.

\begin{alertblock}{Правило}
Если произведение матриц не определено или результат имеет неверную размерность, формула производной записана неправильно.
\end{alertblock}
\end{frame}

% 11
\begin{frame}{Аффинное преобразование: $z=Ax+b$}
\small
Возьмём $A\in\R^{m\times d}$, $x\in\R^d$. Тогда
\[
z(x+dx)-z(x)=A(x+dx)+b-(Ax+b)=A\,dx.
\]
Следовательно,
\[
\boxed{dz=A\,dx,\qquad J_z(x)=A.}
\]
Матрица $A$ и есть линейная часть изменения $z$. Никаких координатных сумм не потребовалось.
\end{frame}

% 12
\begin{frame}{Квадратичная потеря: разделим вычисление на операции}
\small
Пусть
\[
\boxed{f(x)=\frac12\norm{Ax-b}_2^2.}
\]
Введём промежуточный остаток $r=Ax-b$, тогда
\[
 f(x)=\frac12r^\top r.
\]
Теперь достаточно продифференцировать \textbf{две простые операции}:
\[
 x\longmapsto r,\qquad r\longmapsto \tfrac12r^\top r.
\]
\end{frame}

% 13
\begin{frame}{Дифференциал квадрата нормы}
\small
Для $r\in\R^m$ раскрываем квадрат приращения:
\[
\begin{aligned}
\tfrac12(r+dr)^\top(r+dr)
&=\tfrac12r^\top r+r^\top dr+\tfrac12\norm{dr}^2.
\end{aligned}
\]
Линейная часть равна
\[
\boxed{df=r^\top dr.}
\]
Для $r=Ax-b$ уже известно
\[
\boxed{dr=A\,dx.}
\]
\end{frame}

% 14
\begin{frame}{Первое доказательство: градиент квадратичной потери}
\small
\begin{block}{Утверждение}
Для $f(x)=\frac12\norm{Ax-b}_2^2$ выполнено
\[
\boxed{\grad f(x)=A^\top(Ax-b).}
\]
\end{block}
\textbf{Доказательство:}
\[
\begin{aligned}
df&=r^\top dr=r^\top A\,dx\\
  &=(A^\top r)^\top dx.
\end{aligned}
\]
Сопоставляя с $df=\grad f(x)^\top dx$, получаем требуемое. $\square$
\end{frame}

% 16
\begin{frame}{От ошибки одного объекта — к эмпирическому риску}
\small
Даны наблюдения $(x_i,y_i)$, $i=1,\ldots,n$, и линейный прогноз $\hat y_i=w^\top x_i$.
Для одного объекта возьмём квадратичную потерю
\[
\ell_i(w)=\frac12(w^\top x_i-y_i)^2.
\]
\begin{block}{Эмпирический риск}
Средняя потеря \emph{на наблюдённой выборке}:
\[
\boxed{\widehat R_n(w)=\frac1n\sum_{i=1}^n\ell_i(w).}
\]
\end{block}
В отличие от теоретического риска $R(w)=\mathbb E\,\ell(w;(X,Y))$ — среднего по неизвестному распределению данных — здесь усредняем только имеющиеся наблюдения.
\end{frame}
% 17
\begin{frame}{Как отдельные ошибки влияют на каждый параметр}
\small
Обозначим ошибку предсказания на объекте $i$ через
\[
r_i=w^\top x_i-y_i.
\]
Тогда градиент отдельной потери:
\[
\grad\ell_i(w)=r_i x_i.
\]
И градиент эмпирического риска — среднее \textbf{векторов}:
\[
\boxed{\grad\widehat R_n(w)=\frac1n\sum_{i=1}^n r_i x_i.}
\]
Каждая ошибка $r_i$ умножается на вектор признаков $x_i$: так она даёт вклад в \textbf{каждую координату} градиента.
\end{frame}
% 18
\begin{frame}{Матричная запись градиента эмпирического риска}
\small
В матрице $X\in\R^{n\times d}$ строки — это векторы признаков $x_i^\top$. Остатки собираются в столбец
\[
r=Xw-y\in\R^n.
\]
Поэтому средний риск и его градиент записываются компактно:
\[
\widehat R_n(w)=\frac{1}{2n}\norm{Xw-y}_2^2,
\qquad
\boxed{\grad\widehat R_n(w)=\frac1nX^\top r.}
\]
{\centering
\pic[.91]{least_squares_pipeline.pdf}\par}
\vspace{.35em}
{\raggedright\footnotesize
$X^\top$ переводит $n$ ошибок объектов в $d$ компонент градиента; $1/n$ усредняет по выборке.\par}
\end{frame}
% 19
\begin{frame}{Добавим регуляризацию}
\small
Рассмотрим
\[
F_\lambda(w)=\frac1{2n}\norm{Xw-y}^2+\frac\lambda2\norm{w}^2.
\]
Из линейности дифференциала:
\[
\boxed{\grad F_\lambda(w)=\frac1nX^\top(Xw-y)+\lambda w.}
\]
Вклад штрафа прибавляется к уже найденному градиенту основной потери.
\end{frame}

% 20
\begin{frame}{Логистическая потеря: производная по одному аргументу}
\small
Пусть
\[
\ell(z)=\log(1+e^{-z}).
\]
Обычное правило цепочки:
\[
\ell'(z)=\frac{-e^{-z}}{1+e^{-z}}
=\boxed{-\frac1{1+e^z}}.
\]
Это локальная производная \textbf{по скалярному} аргументу $z$.
\end{frame}

% 21
\begin{frame}{Логистическая потеря: откуда берётся аргумент}
\small
Для наблюдения $(x_i,y_i)$ с $y_i\in\{-1,1\}$:
\[
 z=y_iw^\top x_i,
\quad \ell_i(w)=\log(1+e^{-z}).
\]
Функция устроена как композиция:
\[
 w\longmapsto z\longmapsto\ell(z).
\]
Поскольку $z$ линейна по $w$,
\[
\boxed{\grad_w z=y_ix_i.}
\]
\end{frame}

% 22
\begin{frame}{Градиент логистической потери}
\small
По правилу цепочки:
\[
\grad_w \ell_i(w)=\ell'(z)\grad_w z.
\]
Подставим обе уже найденные производные:
\[
\boxed{
\grad_w\ell_i(w)
=-\frac{y_ix_i}{1+e^{y_iw^\top x_i}}.
}
\]
Никакой новой сложной техники здесь нет: только композиция двух функций.
\end{frame}

% 23
\begin{frame}{Общая структура: от параметров до потери}
\small
В рассмотренных моделях вычисление имеет вид
\[
\boxed{\text{параметры}}
\ \longrightarrow\
\boxed{\text{прогноз}}
\ \longrightarrow\
\boxed{\text{потеря}}
\ \longrightarrow\
\boxed{\text{средний риск}}.
\]
Чтобы продифференцировать весь эмпирический риск, достаточно знать производную \textbf{каждой отдельной операции} и правило их композиции.
\begin{alertblock}{Следующая цель}
Найти универсальное правило для отображений между пространствами \emph{разных размерностей}.
\end{alertblock}
\end{frame}
% 24
\begin{frame}{Привычное правило цепочки}
\small
Для одномерных функций:
\[
\boxed{(h\circ g)'(x)=h'(g(x))g'(x).}
\]
Например, для $f(x,y)=\sin(xy)$ вначале вычисляем $g(x,y)=xy$, затем $h(z)=\sin z$.

\begin{block}{Вопрос}
А если промежуточный результат — \textbf{вектор}, а не число?
\end{block}
\end{frame}

% 25
\begin{frame}{Один параметр может влиять по двум путям}
\small
Введём
\[
 a=xy,\qquad b=\sin x,\qquad f=a+b.
\]
Переменная $x$ влияет на $f$ и через $a$, и через $b$:
\[
\frac{\partial f}{\partial x}
=\frac{\partial f}{\partial a}\frac{\partial a}{\partial x}
+\frac{\partial f}{\partial b}\frac{\partial b}{\partial x}
=y+\cos x.
\]
Эти два вклада необходимо \textbf{сложить}.
\end{frame}

% 26
\begin{frame}{Определение матрицы Якоби}
\small
\begin{block}{Определение}
Для дифференцируемой $g:\R^d\to\R^m$ её матрица Якоби — матрица всех первых частных производных:
\[
\boxed{
J_g(x)=
\begin{pmatrix}
\frac{\partial g_1}{\partial x_1}&\cdots&\frac{\partial g_1}{\partial x_d}\\
\vdots&\ddots&\vdots\\
\frac{\partial g_m}{\partial x_1}&\cdots&\frac{\partial g_m}{\partial x_d}
\end{pmatrix}.}
\]
\end{block}
У матрицы $m$ строк (по выходам) и $d$ столбцов (по входам).
\end{frame}

% 27
\begin{frame}{Пример вычисления матрицы Якоби}
\small
Пусть
\[
g(x,y)=\begin{pmatrix}xy\\\sin x\end{pmatrix}.
\]
Производные первой компоненты: $y$, $x$; второй: $\cos x$, $0$. Итак,
\[
\boxed{J_g(x,y)=\begin{pmatrix}y&x\\\cos x&0\end{pmatrix}.}
\]
Строка отвечает за одну \textbf{выходную} координату.
\end{frame}

% 28
\begin{frame}{Теорема: многомерное правило цепочки}
\small
\begin{block}{Теорема}
Пусть $g:\R^d\to\R^m$ дифференцируема в $x$, а $h:\R^m\to\R^p$ дифференцируема в $g(x)$. Тогда
\[
\boxed{J_{h\circ g}(x)=J_h(g(x))\,J_g(x).}
\]
\end{block}
\centering
\pic[.79]{chain_dimensions.pdf}
\end{frame}

% 29
\begin{frame}{Доказательство правила цепочки: два линейных приближения}
\small
Запишем композицию как $x\mapsto z=g(x)\mapsto h(z)$.
Если аргумент меняется на $\Delta x$, промежуточное значение меняется на
\[
\Delta z=g(x+\Delta x)-g(x)
=J_g(x)\Delta x+r_g(\Delta x),
\]
где $r_g(\Delta x)=o(\norm{\Delta x})$.
Для второй функции:
\[
h(z+\Delta z)-h(z)
=J_h(z)\Delta z+r_h(\Delta z),
\]
где $r_h(\Delta z)=o(\norm{\Delta z})$.
\begin{block}{Следующий шаг}
Подставим первое приближение в другое и соберём коэффициент при $\Delta x$.
\end{block}
\end{frame}

% 29a
\begin{frame}{Доказательство правила цепочки: подстановка}
\footnotesize
Подставим $\Delta z=J_g(x)\Delta x+r_g(\Delta x)$:
\[
\begin{aligned}
\Delta(h\circ g)
&=J_h(g(x))J_g(x)\Delta x\\
&\quad+J_h(g(x))r_g(\Delta x)+r_h(\Delta z).
\end{aligned}
\]
Поскольку $\Delta z=O(\norm{\Delta x})$, оба остаточных члена составляют $o(\norm{\Delta x})$. Значит,
\[
\boxed{J_{h\circ g}(x)=J_h(g(x))J_g(x).}
\]
\begin{alertblock}{Проверка размерностей}
$(p\times m)(m\times d)=p\times d$: промежуточное пространство $\R^m$ «сокращается».
\end{alertblock}
\end{frame}
% 30
\begin{frame}{Для скалярного выхода правило становится проще}
\small
Если $h:\R^m\to\R$, то $J_h=\grad h^\top$ — строка. Получаем
\[
J_{h\circ g}(x)=\grad h(g(x))^\top J_g(x).
\]
Транспонируем, чтобы получить градиент-столбец:
\[
\boxed{\grad(h\circ g)(x)=J_g(x)^\top\grad h(g(x)).}
\]
\textbf{Этот порядок матричного произведения} ляжет в основу обратного режима автоматического дифференцирования.
\end{frame}

% 31
\begin{frame}{Применение правила цепочки к знакомой функции}
\small
Пусть
\[
g(x,y)=\begin{pmatrix}xy\\\sin x\end{pmatrix},\qquad
h(u,v)=(u+v)^2.
\]
Тогда
\[
\grad h(u,v)=2(u+v)\begin{pmatrix}1\\1\end{pmatrix},
\quad
J_g(x,y)^\top=\begin{pmatrix}y&\cos x\\x&0\end{pmatrix}.
\]
Следовательно,
\[
\boxed{\grad(h\circ g)(x,y)=2(xy+\sin x)
\begin{pmatrix}y+\cos x\\x\end{pmatrix}.}
\]
\end{frame}

% 33
\begin{frame}{Зададим сложную функцию как последовательность операций}
\small
Будем исследовать одну функцию до конца лекции:
\[
\boxed{f(x,y)=(xy+\sin x)^2.}
\]
Введём промежуточные значения
\[
a=xy,\qquad b=\sin x,\qquad c=a+b,\qquad f=c^2.
\]
Вычислительный граф хранит \textbf{зависимости} между этими операциями.
\end{frame}

% 34
\begin{frame}{Вычислительный граф}
\centering
\pic[.93]{graph_overview.pdf}

\small
Узлы — результаты элементарных операций, стрелки — зависимости. Порядок вычислений определяется направлением стрелок.
\end{frame}

% 35
\begin{frame}{Прямой проход: начальные значения}
\small
Пусть $x=1$, $y=2$. Сначала значения известны только во входных узлах.
\centering
\pic[.90]{graph_forward_start.pdf}
\end{frame}

% 36
\begin{frame}{Прямой проход: промежуточные вычисления}
\small
Следующий уровень можно вычислить независимо:
\[
a=xy=2,\qquad b=\sin x=\sin1.
\]
\centering
\pic[.82]{graph_forward_middle.pdf}
\end{frame}

% 37
\begin{frame}{Прямой проход: вычисляем конечный результат}
\small
Теперь $c=2+\sin1$, затем $f=(2+\sin1)^2$.
\centering
\pic[.90]{graph_forward_values.pdf}
\end{frame}

% 38
\begin{frame}{Локальные производные в каждом узле}
\small
Почти все необходимые производные элементарны:
\[
\begin{array}{ll}
a=xy:& a_x=y,\quad a_y=x,\\[2pt]
b=\sin x:& b_x=\cos x,\\[2pt]
c=a+b:& c_a=c_b=1,\\[2pt]
f=c^2:& f_c=2c.
\end{array}
\]
\begin{block}{Наблюдение}
Вычислить \emph{одну локальную производную} гораздо проще, чем сразу получить градиент всей длинной композиции.
\end{block}
\end{frame}

% 39
\begin{frame}{Два пути влияния одного параметра}
\centering
\pic[.92]{two_paths.pdf}

\small
\[
\boxed{\frac{\partial f}{\partial x}=2c(y+\cos x),\qquad
\frac{\partial f}{\partial y}=2cx.}
\]
\end{frame}

% 40
\begin{frame}{Зачем нам два способа проходить по графу?}
\small
Пусть у функции $d$ входных параметров, а на выходе $m$ чисел:
\[
F:\R^d\longrightarrow\R^m.
\]
Нужно уметь считать, \textbf{как малое изменение входов отражается на выходах}.
\begin{itemize}
\item В \textbf{прямом режиме} выбираем изменение входов $v$ и отслеживаем, как оно проходит через операции. Получаем $J_F(x)v$.
\item В \textbf{обратном режиме} выбираем скалярную комбинацию выходов $u^\top F$ и выясняем, \emph{через какие промежуточные узлы} она зависит от входов. Получаем $J_F(x)^\top u$.
\end{itemize}
\begin{alertblock}{Что такое «вклад»?}
Если вход влияет на результат двумя путями, его производная равна \textbf{сумме вкладов} каждого пути по правилу цепочки.
\end{alertblock}
\end{frame}

% 40a
\begin{frame}{Два режима: что распространяется по графу?}
{\centering
\pic[.86]{forward_reverse_compare.pdf}\par}
\vspace{.25em}
{\raggedright\small
\textbf{Прямой проход:} от входов к выходу — вычисляем значения и направленные производные.\par
\vspace{.35em}
\textbf{Обратный проход:} после вычисления значений идём от выхода к входам и накапливаем производные.\par}
\end{frame}
% 41
\begin{frame}{Прямой режим: что означает точка над переменной?}
\small
Вообразим, что исходные аргументы меняются с параметром $t$:
\[
x(t)=x+t v_x,\qquad y(t)=y+t v_y.
\]
Выбранное направление — это вектор $v=(v_x,v_y)$. \textbf{Точкой} обозначаем скорость изменения узла при $t=0$:
\[
\boxed{\dot u:=\left.\frac{du(t)}{dt}\right|_{t=0}.}
\]
Например, для $a(t)=x(t)y(t)$:
\[
\dot x=v_x,\quad\dot y=v_y,
\qquad\boxed{\dot a=y\dot x+x\dot y.}
\]
Если хотим $\partial f/\partial x$, берём $v=(1,0)$: меняется только $x$, а $y$ остаётся постоянным.
\end{frame}
% 42
\begin{frame}{Прямой режим: локальные правила}
\small
\[
\begin{array}{lll}
a=u+v &\Longrightarrow& \dot a=\dot u+\dot v,\\[7pt]
a=uv &\Longrightarrow& \dot a=v\dot u+u\dot v,\\[7pt]
a=\sin u &\Longrightarrow& \dot a=\cos(u)\dot u,\\[7pt]
a=u^2 &\Longrightarrow& \dot a=2u\dot u.
\end{array}
\]
Это ровно знакомые правила дифференцирования, применяемые при обходе графа.
\end{frame}

% 43
\begin{frame}{Прямой режим: находим $\partial f/\partial x$}
\small
Начинаем с $\dot x=1$, $\dot y=0$:
\[
\dot a=y\dot x+x\dot y=y,\qquad
\dot b=\cos x\,\dot x=\cos x.
\]
\[
\dot c=\dot a+\dot b=y+\cos x.
\]
В конечном узле:
\[
\boxed{\dot f=2c\dot c=2(xy+\sin x)(y+\cos x).}
\]
\end{frame}

% 44
\begin{frame}{Прямой режим: как найти другую компоненту?}
\small
Для $\partial f/\partial y$ изменяем только начальные производные:
\[
\boxed{\dot x=0,\qquad \dot y=1.}
\]
Тогда
\[
\dot a=x,\quad \dot b=0,\quad\dot c=x,
\quad\boxed{\dot f=2cx.}
\]
\begin{alertblock}{Ограничение прямого режима}
Для полного градиента скалярной функции $d$ переменных в общем случае нужны $d$ независимых направленных проходов.
\end{alertblock}
\end{frame}

% 45
\begin{frame}{Обратный режим: поменяем направление обхода}
\small
Сначала выполняем \textbf{обычный прямой проход}, чтобы знать значения всех узлов.

Затем начинаем с конечной функции. Вводим обозначение:
\[
\boxed{\bar u:=\frac{\partial f}{\partial u}.}
\]
В конечном узле всегда $\bar f=1$. Остальные величины получим, проходя от выхода к входам.
\end{frame}

% 46
\begin{frame}{Обратный проход: от результата к сумме}
\centering
\pic[.78]{graph_reverse_seed.pdf}

\small
Так как $f=c^2$ и $\bar f=1$, получаем
\[
\boxed{\bar c=2c.}
\]
На схеме стрелки показывают порядок \textbf{распространения производных}, а не вычисления значений.
\end{frame}

% 47
\begin{frame}{Обратный проход: разветвление и накопление}
\small
Поскольку $c=a+b$, от узла $c$ получаем
\[
\bar a=\bar c\frac{\partial c}{\partial a}=2c,
\qquad
\bar b=\bar c\frac{\partial c}{\partial b}=2c.
\]
\centering
\pic[.69]{graph_reverse_mid.pdf}
\end{frame}

% 48
\begin{frame}{Обратный проход: собираем градиент по входам}
\small
От ветки $a=xy$:
\[
\bar x\mathrel{+}=\bar a\,y,
\qquad\bar y\mathrel{+}=\bar a\,x.
\]
От ветки $b=\sin x$:
\[
\bar x\mathrel{+}=\bar b\cos x.
\]
В итоге
\[
\boxed{\bar x=2c(y+\cos x),\qquad \bar y=2cx.}
\]
\begin{alertblock}{Ключевой принцип}
Если в параметр ведёт несколько путей, их вклады в производную \textbf{складываются}.
\end{alertblock}
\end{frame}

% 49
\begin{frame}{Какой режим выгоднее?}
\small
\begin{center}
\begin{tabular}{p{.43\linewidth}p{.43\linewidth}}
\toprule
\textbf{Прямой режим} & \textbf{Обратный режим}\\
\midrule
Вычисляет $Jv$ & Вычисляет $J^\top u$\\
\addlinespace
Особенно удобен при небольшом числе входов & Особенно удобен при небольшом числе выходов\\
\addlinespace
Для полного градиента нужен проход по каждому базисному направлению & Для скалярного выхода достаточно одного обратного прохода\\
\bottomrule
\end{tabular}
\end{center}
\smallskip
\textbf{Цена обратного режима:} часто нужно сохранять промежуточные значения (или пересчитывать их), то есть возникает компромисс между памятью и временем.
\end{frame}

% 50
\begin{frame}{А можно просто приближать производные числами?}
\small
Например, центральная конечная разность:
\[
\boxed{f'(x)\approx\frac{f(x+h)-f(x-h)}{2h}.}
\]
При достаточно гладкой функции погрешность усечения имеет порядок $O(h^2)$, но слишком малый шаг усиливает влияние ошибок округления.

\begin{alertblock}{Вопрос}
Почему нельзя просто выбрать $h=10^{-100}$ и получить идеальную производную?
\end{alertblock}
\end{frame}

% 51
\begin{frame}{Численная производная: большой и очень маленький шаг}
\centering
\pic[.81]{finite_difference_error.pdf}

\small
При слишком большом $h$ приближение грубо; при слишком маленьком возникают потери точности из-за вычитания близких чисел.
\end{frame}

% 52
\begin{frame}[fragile]{Как выглядит автоматическое дифференцирование в Python}
\footnotesize
\begin{semiverbatim}
import torch
x = torch.tensor(1.0, requires_grad=True)
y = torch.tensor(2.0, requires_grad=True)
f = (x*y + torch.sin(x))**2
f.backward()
print(x.grad.item(), y.grad.item())
\end{semiverbatim}
\small
Команда \texttt{backward()} выполняет обратный проход и заполняет производные у входных переменных.

\begin{block}{Математика не изменилась}
Компьютер применяет то же правило цепочки, которое мы только что проделали вручную.
\end{block}
\end{frame}

% 53
\begin{frame}{Проверяем три способа вычислить производную}
\small
Для $f(x,y)=(xy+\sin x)^2$ можно сравнить:
\begin{itemize}
\item \textbf{Аналитическую:} $\partial_xf=2(xy+\sin x)(y+\cos x)$;
\item \textbf{Численную:} конечную разность по $x$;
\item \textbf{Автоматическую:} результат обратного прохода.
\end{itemize}
\begin{alertblock}{Как интерпретировать согласие}
Совпадение результатов — полезная проверка реализации, но не формальное доказательство правильности для всех входов.
\end{alertblock}
\end{frame}

% 54
\begin{frame}{Три частые ошибки в реальной реализации}
\small
\begin{enumerate}
\item \textbf{Размерности:} записали $A(Ax-b)$ вместо $A^\top(Ax-b)$.
\item \textbf{Масштаб потери:} при переходе от суммы к среднему забыли множитель $1/n$.
\item \textbf{Накопление:} повторно вызвали обратный проход, не очистив ранее накопленные градиенты в библиотеке.
\end{enumerate}
\begin{exampleblock}{Почему помогает теория}
Если понимать, что именно вычисляет градиент, такие ошибки легче заметить, чем при бездумном использовании библиотеки.
\end{exampleblock}
\end{frame}

% 55
\begin{frame}{Карта лекции}
\small
\[
\boxed{df=J_f\,dx}
\]
\[
\Downarrow\quad
\boxed{J_{h\circ g}=J_hJ_g}
\]
\[
\Downarrow\quad
\boxed{\text{локальные производные на графе}}
\]
\[
\Downarrow\quad
\boxed{\text{прямой и обратный проходы}}
\]
Автоматическое дифференцирование — алгоритмическое применение правила цепочки, а не численная аппроксимация конечными разностями.
\end{frame}

% 56
\begin{frame}{Следующая лекция: как выбирать шаг?}
\small
Теперь мы знаем, \textbf{как вычислить} $\grad f(x_k)$ даже для сложной функции.

Но градиентный спуск требует ещё одного решения:
\[
\boxed{x_{k+1}=x_k-\alpha_k\grad f(x_k).}
\]
\begin{alertblock}{Следующий вопрос}
Как выбирать $\alpha_k$ и когда останавливать алгоритм, если оптимальное решение заранее неизвестно?
\end{alertblock}
Далее — поиск шага, достаточное убывание и критерии остановки.
\end{frame}

\end{document}
